\chapter*{ABSTRACT}
\addcontentsline{toc}{chapter}{ABSTRACT}

Manual food logging often fails for a simple reason: the same work has to be repeated after every meal. A user has to identify each dish, estimate the quantity, search for a suitable database entry, and then repeat the process for the next item on the plate. This becomes more difficult for Indian meals because one plate can contain rice, dal, chutney, curry, fried snacks, sweets, and beverages, sometimes with regional or household-specific names. This project addresses that practical problem by developing an AI-based food recognition system with nutrition-aware recommendations. The system accepts a food image, detects visible food items using a YOLO-based object detector, maps the detected labels to Indian nutrition records, estimates calories and macronutrients per 100 grams, and presents goal-aware dietary suggestions through a web application.

The implemented prototype combines a FastAPI backend, a Next.js frontend, a SQLite nutrition database built mainly from the Indian Nutrient Database, and a merged YOLO-format dataset containing 72 canonical Indian food classes. Five source datasets were consolidated through a normalization pipeline that resolves spelling variants, regional labels, and duplicate class names. This step required careful checking because small naming differences such as \code{idly} and \code{idli}, or \code{satham} and \code{rice}, can become separate detector classes if they are not handled before training. The final balanced dataset contains 48,693 exported image-label pairs across training, validation, and test splits. Visual verification sheets were generated from real dataset images, confirming that all 72 class names have corresponding samples while also revealing a small number of annotation issues that should be corrected in later training cycles.

The final detector training was completed with YOLO11s on the 72-class dataset. The best model achieved 0.830 mAP@50 and 0.692 mAP@50:95 on validation, and 0.820 mAP@50 and 0.680 mAP@50:95 on the held-out test split. The nutrition layer maps all 72 dataset classes to database entries by using explicit semantic mappings and verified supplemental rows before fuzzy matching. This design decision was taken after observing that simple fuzzy matching can return a food with a similar name but a different nutritional profile. The backend exposes endpoints for image analysis, nutrition validation, macro calculation, daily goals, health conditions, and recommendation generation. The frontend supports image upload, detection visualization, serving adjustment, macro tracking, and personalized guidance. The report presents the design decisions, implementation process, validation evidence, limitations, and possible future improvements of the project.
